\section{Proofs of Estimators (1D GMM - Classification)}

For $p=1$, there is only one feature, so we consider a
\hl{one-dimensional Gaussian mixture model}:
\begin{equation*}
X\mid Z=k\sim\mathcal{N}(\mu_k,\sigma_k^2),
\qquad k=1,\ldots,K.
\end{equation*}

The one-hot representation of the group variable is
\begin{equation*}
\tilde{Z}\sim\mathcal{M}(1,p_1,\ldots,p_K),
\end{equation*}
where
\begin{equation*}
p_k=P(Z=k),
\qquad
\sum_{k=1}^Kp_k=1.
\end{equation*}

The parameter vector is
\begin{equation*}
\theta=
(p_1,\ldots,p_K,\mu_1,\ldots,\mu_K,
\sigma_1^2,\ldots,\sigma_K^2).
\end{equation*}

For group $k$, the Gaussian density is
\begin{equation*}
f_k(x)
=
\frac{1}{\sqrt{2\pi\sigma_k^2}}
\exp\left(
-\frac{1}{2}
\frac{(x-\mu_k)^2}{\sigma_k^2}
\right).
\end{equation*}

We observe
\begin{equation*}
\underline{X}=(X_1,\ldots,X_n)
\end{equation*}
and denote the corresponding group variables by
\begin{equation*}
\underline{Z}=(Z_1,\ldots,Z_n).
\end{equation*}

Using the one-hot representation $\tilde{Z}_{ik}$, where
\begin{equation*}
\tilde{Z}_{ik}
=
\begin{cases}
1 & \text{if } Z_i=k,\\
0 & \text{otherwise},
\end{cases}
\end{equation*}
the complete-data log-likelihood is
\begin{equation*}
\ln L(\underline{X},Z,\theta)
=
\sum_{i=1}^{n}\sum_{k=1}^{K}
\tilde{Z}_{ik}
\left[
\ln p_k
-\frac{1}{2}\ln\sigma_k^2
-\ln\sqrt{2\pi}
-\frac{1}{2}
\frac{(x_i-\mu_k)^2}{\sigma_k^2}
\right].
\end{equation*}

We now derive the maximum likelihood estimators of the three types
of parameters:
\begin{equation*}
p_k,\qquad \mu_k,\qquad \sigma_k^2.
\end{equation*}

\subsection{Estimation of the Mixing Proportions $p_k$}

The mixing proportions satisfy the constraint
\begin{equation*}
\sum_{k=1}^Kp_k=1.
\end{equation*}

Equivalently,
\begin{equation*}
\sum_{k=1}^Kp_k-1=0.
\end{equation*}

Following the notation used by the professor, we can write
\begin{equation*}
\sum_{k=1}^Kp_k-1=\psi(x).
\end{equation*}

To find the extremum of a function $\phi(x)$ under the constraint
\begin{equation*}
\psi(x)=0,
\end{equation*}
we introduce a Lagrange multiplier $\lambda$ and optimize
\begin{equation*}
F(x,\lambda)
=
\phi(x)+\lambda\psi(x).
\end{equation*}

In our case, we therefore define the Lagrangian
\begin{equation*}
F(\theta,\lambda)
=
\ln L(\underline{X},Z,\theta)
+
\lambda
\left(
1-\sum_{k=1}^Kp_k
\right).
\end{equation*}

We use the conditions
\begin{equation*}
\frac{\partial F}{\partial p_k}=0,
\qquad
\frac{\partial F}{\partial\lambda}=0.
\end{equation*}

When differentiating with respect to $p_k$, all terms that do not
depend on $p_k$ can be ignored. Therefore,
\begin{equation*}
F
=
\sum_{k=1}^K
\left(
\sum_{i=1}^n\tilde{Z}_{ik}
\right)\ln p_k
+
\lambda
\left(
1-\sum_{k=1}^Kp_k
\right)
+
\text{terms independent of the }p_k.
\end{equation*}

Define
\begin{equation*}
n_k
=
\sum_{i=1}^n\tilde{Z}_{ik},
\end{equation*}
where $n_k$ is the number of observations belonging to group $k$.

Thus, the part of the Lagrangian that depends on the mixing
proportions becomes
\begin{equation*}
F
=
\sum_{k=1}^K n_k\ln p_k
+
\lambda
\left(
1-\sum_{k=1}^Kp_k
\right)
+
\text{constant}.
\end{equation*}

\paragraph{Step 1: Differentiate with respect to $p_k$.}

For a fixed group $k$,
\begin{equation*}
\frac{\partial F}{\partial p_k}
=
\frac{n_k}{p_k}-\lambda.
\end{equation*}

At the optimum,
\begin{equation*}
\frac{n_k}{p_k}-\lambda=0.
\end{equation*}

Therefore,
\begin{equation*}
p_k=\frac{n_k}{\lambda}.
\end{equation*}

\paragraph{Step 2: Differentiate with respect to $\lambda$.}

The derivative with respect to the Lagrange multiplier gives back
the constraint:
\begin{equation*}
\frac{\partial F}{\partial\lambda}
=
1-\sum_{k=1}^Kp_k
=
0.
\end{equation*}

Therefore,
\begin{equation*}
\sum_{k=1}^Kp_k=1.
\end{equation*}

Substituting
\begin{equation*}
p_k=\frac{n_k}{\lambda}
\end{equation*}
into the constraint gives
\begin{equation*}
\sum_{k=1}^K\frac{n_k}{\lambda}=1.
\end{equation*}

Since every observation belongs to exactly one group,
\begin{equation*}
\sum_{k=1}^Kn_k=n.
\end{equation*}

Therefore,
\begin{equation*}
\frac{n}{\lambda}=1
\qquad\Longrightarrow\qquad
\lambda=n.
\end{equation*}

Hence,
\begin{equation*}
\boxed{
\hat{p}_k=\frac{n_k}{n}
}.
\end{equation*}

Since
\begin{equation*}
n_k=\sum_{i=1}^n\tilde{Z}_{ik},
\end{equation*}
we can equivalently write
\begin{equation*}
\boxed{
\hat{p}_k
=
\frac{1}{n}
\sum_{i=1}^n\tilde{Z}_{ik}
}.
\end{equation*}

Thus, the estimated mixing proportion of group $k$ is simply the
\hl{proportion of observations belonging to group $k$}.

\subsection{Estimation of the Mean $\mu_k$}

We now estimate the mean $\mu_k$ of group $k$.

Unlike the mixing proportions, the means are not subject to a
constraint. Therefore, we can directly differentiate the
log-likelihood with respect to $\mu_k$.

Starting from
\begin{equation*}
\ln L(\underline{X},Z,\theta)
=
\sum_{i=1}^{n}\sum_{k=1}^{K}
\tilde{Z}_{ik}
\left[
\ln p_k
-\frac{1}{2}\ln\sigma_k^2
-\ln\sqrt{2\pi}
-\frac{1}{2}
\frac{(x_i-\mu_k)^2}{\sigma_k^2}
\right],
\end{equation*}
we only keep the terms that depend on $\mu_k$:
\begin{equation*}
\ln L
=
-\frac{1}{2}
\sum_{i=1}^n
\tilde{Z}_{ik}
\frac{(x_i-\mu_k)^2}{\sigma_k^2}
+
\text{constant}.
\end{equation*}

We differentiate with respect to $\mu_k$:
\begin{equation*}
\frac{\partial\ln L}{\partial\mu_k}
=
-\frac{1}{2}
\sum_{i=1}^n
\tilde{Z}_{ik}
\frac{\partial}{\partial\mu_k}
\left(
\frac{(x_i-\mu_k)^2}{\sigma_k^2}
\right).
\end{equation*}

Since $\sigma_k^2$ does not depend on $\mu_k$,
\begin{equation*}
\frac{\partial}{\partial\mu_k}
(x_i-\mu_k)^2
=
-2(x_i-\mu_k).
\end{equation*}

Therefore,
\begin{equation*}
\frac{\partial\ln L}{\partial\mu_k}
=
\frac{1}{\sigma_k^2}
\sum_{i=1}^n
\tilde{Z}_{ik}(x_i-\mu_k).
\end{equation*}

At the maximum likelihood estimator,
\begin{equation*}
\frac{\partial\ln L}{\partial\mu_k}=0.
\end{equation*}

Thus,
\begin{equation*}
\frac{1}{\sigma_k^2}
\sum_{i=1}^n
\tilde{Z}_{ik}(x_i-\mu_k)
=
0.
\end{equation*}

Since $\sigma_k^2>0$, we can multiply by $\sigma_k^2$:
\begin{equation*}
\sum_{i=1}^n
\tilde{Z}_{ik}(x_i-\mu_k)
=
0.
\end{equation*}

Expanding the expression,
\begin{equation*}
\sum_{i=1}^n
\tilde{Z}_{ik}x_i
-
\mu_k
\sum_{i=1}^n
\tilde{Z}_{ik}
=
0.
\end{equation*}

Using
\begin{equation*}
n_k=\sum_{i=1}^n\tilde{Z}_{ik},
\end{equation*}
we obtain
\begin{equation*}
\sum_{i=1}^n
\tilde{Z}_{ik}x_i
-
\mu_kn_k
=
0.
\end{equation*}

Therefore,
\begin{equation*}
\mu_kn_k
=
\sum_{i=1}^n
\tilde{Z}_{ik}x_i
\end{equation*}
and hence
\begin{equation*}
\boxed{
\hat{\mu}_k
=
\frac{1}{n_k}
\sum_{i=1}^n
\tilde{Z}_{ik}x_i
}.
\end{equation*}

Thus, the maximum likelihood estimator of $\mu_k$ is simply the
\hl{sample mean of the observations belonging to group $k$}.

\subsection{Estimation of the Variance $\sigma_k^2$}

Finally, we estimate the variance $\sigma_k^2$ of group $k$.

Again, there is no constraint on $\sigma_k^2$, so we directly
differentiate the log-likelihood with respect to $\sigma_k^2$.

The terms depending on $\sigma_k^2$ are
\begin{equation*}
\ln L
=
\sum_{i=1}^n
\tilde{Z}_{ik}
\left[
-\frac{1}{2}\ln\sigma_k^2
-\frac{1}{2}
\frac{(x_i-\mu_k)^2}{\sigma_k^2}
\right]
+
\text{constant}.
\end{equation*}

Therefore,
\begin{equation*}
\frac{\partial\ln L}{\partial\sigma_k^2}
=
\sum_{i=1}^n
\tilde{Z}_{ik}
\frac{\partial}{\partial\sigma_k^2}
\left[
-\frac{1}{2}\ln\sigma_k^2
-\frac{1}{2}
\frac{(x_i-\mu_k)^2}{\sigma_k^2}
\right].
\end{equation*}

For the first term,
\begin{equation*}
\frac{\partial}{\partial\sigma_k^2}
\left(
-\frac{1}{2}\ln\sigma_k^2
\right)
=
-\frac{1}{2\sigma_k^2}.
\end{equation*}

For the second term,
\begin{equation*}
\frac{\partial}{\partial\sigma_k^2}
\left(
-\frac{1}{2}
\frac{(x_i-\mu_k)^2}{\sigma_k^2}
\right)
=
\frac{1}{2}
\frac{(x_i-\mu_k)^2}{(\sigma_k^2)^2}.
\end{equation*}

Therefore,
\begin{equation*}
\frac{\partial\ln L}{\partial\sigma_k^2}
=
\sum_{i=1}^n
\tilde{Z}_{ik}
\left[
-\frac{1}{2\sigma_k^2}
+
\frac{1}{2}
\frac{(x_i-\mu_k)^2}{(\sigma_k^2)^2}
\right].
\end{equation*}

At the maximum likelihood estimator,
\begin{equation*}
\frac{\partial\ln L}{\partial\sigma_k^2}=0.
\end{equation*}

Thus,
\begin{equation*}
\sum_{i=1}^n
\tilde{Z}_{ik}
\left[
-\frac{1}{2\sigma_k^2}
+
\frac{1}{2}
\frac{(x_i-\mu_k)^2}{(\sigma_k^2)^2}
\right]
=
0.
\end{equation*}

Multiplying by $2(\sigma_k^2)^2$ gives
\begin{equation*}
\sum_{i=1}^n
\tilde{Z}_{ik}
\left[
-\sigma_k^2
+
(x_i-\mu_k)^2
\right]
=
0.
\end{equation*}

Therefore,
\begin{equation*}
\sum_{i=1}^n
\tilde{Z}_{ik}(x_i-\mu_k)^2
=
\sigma_k^2
\sum_{i=1}^n\tilde{Z}_{ik}.
\end{equation*}

Using
\begin{equation*}
n_k=\sum_{i=1}^n\tilde{Z}_{ik},
\end{equation*}
we obtain
\begin{equation*}
\sum_{i=1}^n
\tilde{Z}_{ik}(x_i-\mu_k)^2
=
n_k\sigma_k^2.
\end{equation*}

Hence,
\begin{equation*}
\boxed{
\hat{\sigma}_k^2
=
\frac{1}{n_k}
\sum_{i=1}^n
\tilde{Z}_{ik}(x_i-\hat{\mu}_k)^2
}.
\end{equation*}

Thus, the maximum likelihood estimator of $\sigma_k^2$ is the
\hl{average squared deviation from the estimated group mean},
calculated using only the observations belonging to group $k$.

\subsubsection*{Summary of the Three Estimators}

For a one-dimensional Gaussian mixture with known group memberships,
the maximum likelihood estimators are
\begin{equation*}
\boxed{
\hat{p}_k
=
\frac{n_k}{n}
}
\end{equation*}
\begin{equation*}
\boxed{
\hat{\mu}_k
=
\frac{1}{n_k}
\sum_{i=1}^n
\tilde{Z}_{ik}x_i
}
\end{equation*}
and
\begin{equation*}
\boxed{
\hat{\sigma}_k^2
=
\frac{1}{n_k}
\sum_{i=1}^n
\tilde{Z}_{ik}(x_i-\hat{\mu}_k)^2
}.
\end{equation*}

In other words,
\begin{equation*}
\boxed{
\begin{aligned}
\hat{p}_k
&=\text{proportion of observations in group }k,\\[2mm]
\hat{\mu}_k
&=\text{mean of the observations in group }k,\\[2mm]
\hat{\sigma}_k^2
&=\text{variance of the observations in group }k.
\end{aligned}
}
\end{equation*}

These are exactly the usual empirical estimators applied separately
within each group.